\section*{Five holographic principles of construction and conservation}
\addcontentsline{toc}{section}{Five holographic principles of construction and conservation}

The \emph{Elements} provide a precise antecedent for presenting a
mathematical architecture through declared operations and proved
consequences. Book I distinguishes definitions, postulates, and common notions.
Its first three postulates allow a segment to be drawn between points,
a segment to be extended in a straight line, and a circle to be described with a
given centre and radius; the fourth establishes the equality of right angles.
The fifth formulates a condition of intersection: two straight lines cut by a transversal
meet, when extended, on the side where the interior angles sum to
less than two right angles. These are statements about the objects and constructions of
that geometry.\footnote{Euclid, \emph{Elements}, Book I, Postulates 1--5.
Heiberg's Greek text and Richard Fitzpatrick's translation, p.~7:
\url{https://farside.ph.utexas.edu/books/Euclid/Elements.pdf}. They may also
be consulted separately in David E. Joyce's edition:
\url{https://mathcs.clarku.edu/~djoyce/elements/bookI/bookI.html}.}

The extreme and mean ratio appears in Book VI as a relation between a
whole and its parts: the whole is to the greater segment as the latter is to the lesser.
In a modern algebraic translation, for positive lengths \(u>v\), one writes
\((u+v)/u=u/v\); hence the ratio \(r=u/v\) satisfies
\(r^2=r+1\). This notation explains the transmitted proportion without turning
modern algebraic language into Euclid's starting point. In HMT,
the self-scaling coordinate is obtained from an incidence produced by the
operational calendar; its recognition through the same equation belongs
to a subsequent stage.\footnote{Euclid, \emph{Elements}, Book VI,
Definition 3, in Joyce's edition:
\url{https://mathcs.clarku.edu/~djoyce/elements/bookVI/defVI3.html}.
Fitzpatrick's edition prints this same statement as Definition 2
of Book VI, p.~156. Both editorial numberings are retained here.}

The comparison that follows concerns the organization of the construction.
HMT declares an arithmetic support, two related sheets, a tritic
orientation, and operations that prolong states while preserving their records.
A cylinder represents the information fixed by a prefix; a Euclidean
line is another object, with another definition and other operations. The
five principles in this section express effective relations between
holographic objects. Their number provides a recognizable exposition; it
does not establish a term-by-term correspondence with the five postulates,
a refutation of the fifth, or a replacement for the geometric axioms.

These principles have an \emph{organizing} status: they bring together conditions and
theorems that the body of the article constructs and proves. Each first presents
its meaning, then its domains and an identity by which it can be checked.
The common sequence is
\[
 \mathrm{APP}\longrightarrow\mathrm{TRIT}\longrightarrow\mathrm{TPK}
 \longrightarrow\text{enriched state}
 \longrightarrow\text{joint discrete structure of the continuum}.
\]
Numerical, linear, and geometric realizations subsequently read that
structure. This placement of the realizations keeps visible the origin
of the coefficients and the information preserved by each reading.

\subsection*{Principle I. Relational unity of generation}
\hypertarget{prinh:unidad}{}

\emph{Each reading is constructed through typed operations on a
common support, preserving the relations that allow it to be continued.} An
APP cell simultaneously admits an additive and a multiplicative reading.
The passage to a visible digit records part of both operations; the
quotient, orientation, and sheet preserve the conditions of its
prolongation. Unity consists in this relational provenance, which allows
the same execution to be followed through its different realizations.

The initial domain is \(D_9^2\), with \(D_9=\{1,\ldots,9\}\). On it
act \(S(i,j)=i+j\) and \(P(i,j)=ij\). For an integer \(n\), the
positive residue-and-quotient chart is given by
\[
 \rho_9(n)=1+((n-1)\bmod9),\qquad
 q_9(n)=\frac{n-\rho_9(n)}9,\qquad
 n=\rho_9(n)+9q_9(n).
\]
The lifted APP evaluation preserves
\(\bigl((\rho_9(S),q_9(S)),(\rho_9(P),q_9(P))\bigr)\).
The identity reconstructs both integer readings; for example, it distinguishes
the values 10 and 19, which have the same positive residue and different
quotients. The result and its proof are developed in
\eqref{eq:residuo-cociente} and \eqref{eq:app-levantada}. The representation
by residues \(0,\ldots,8\), used in other charts of the article, preserves
its corresponding quotient: changing representatives requires changing both
components of the division.

TRIT adds the oriented representation of \(\mathbb F_3\) by
\(\{-1,0,+1\}\), with its exact division \(n=r+3c\). The operational selector
\(M_{\rm ph}\) takes the values \(+1,-1,0\) in the phase classes
\(\{1,4,7\}\), \(\{2,5,8\}\), \(\{3,6,9\}\), respectively.
Cardinal transport has another domain and another function: it moves the
cursors on \(X_9=(\mathbb Z/9\mathbb Z)^2\). In the enriched state
\(\widehat X\), TPK composes selection, transport, and memory
updating,
\[
 U_t=\operatorname{Mem}_t\circ\operatorname{Tra}_t
                    \circ\operatorname{Sel}_t.
\]
Definition \eqref{eq:actualizacion} and the development in
\cref{tpk-int:estado,tpk-int:seleccion} specify phase, cursors, sheets,
directions, residues, quotients, and records. Phase sign and cardinal
direction are composed while preserving their distinct types.

Self-scaling provides an example of generation that can be followed without starting
from its final value. The calendar \((+1,-1,0)\) produces, with the rule for
changing and preserving sheets, the modal cycle
\((\Sigma,\Pi,\Pi)\). Its three edges determine
\[
 F_{\rm av}=\begin{pmatrix}0&1\\1&1\end{pmatrix}.
\]
In its subsequent linear realization, the positive eigenray normalized
as \((1,x)^{\mathsf T}\) satisfies \(x=\lambda\) and
\(1+x=\lambda x\). Thus \(x^2-x-1=0\) is obtained; the positive root is
recognized as \(\varphi\). The theorem following \eqref{eq:fav} proves
this identification after constructing the entries of the matrix.
The relation to the Euclidean proportion is exact in this self-scaling
reading and preserves the distinction between their two procedures of origin.

\subsection*{Principle II. Oriented prolongation and composition memory}
\hypertarget{prinh:prolongacion}{}

\emph{An admissible prolongation adds an event to the state and
transports the preceding memory through the same operation.} The history
allows paths ending at the same position or phase to be
distinguished. This principle provides a precise way of speaking about
continuation: the prefix actually traversed is preserved, and an
emission allowed by the current selector and regime is added.

Let \(X_k\) be the set of admissible histories of length \(k\),
constructed from the seeds of the kernel. If \(\varepsilon\) is admissible
after \(h\), the prolongation \(h\varepsilon\) belongs to \(X_{k+1}\)
and truncation satisfies \(\rho_k(h\varepsilon)=h\). The
emission label and the preceding sheets are retained, even when two emissions
have the same final reading. In a memory chart over a module
\(M\), the update has the form
\(U_\varepsilon(m)=C_\varepsilon m+\kappa_\varepsilon\), with
\(C_\varepsilon\in\operatorname{End}(M)\) and
\(\kappa_\varepsilon\in M\). If two segments are composable in that
chart, then
\[
 C_{\beta\alpha}=C_\beta C_\alpha,\qquad
 \kappa_{\beta\alpha}=C_\beta\kappa_\alpha+\kappa_\beta,
 \qquad U_{\beta\alpha}=U_\beta U_\alpha.
\]
The proof of \cref{iv:cont:concatenacion} consists in substituting the first
update into the second. Memory is added after the transport
corresponding to it; associativity makes the result independent of the
subdivision of the path.

The distinction between phase and memory already appears in an elementary reading.
For \((w,r)\in\mathbb Z\times C_9\), with
\(\bar r\in\{0,\ldots,8\}\), one has
\[
 \sigma_9(w,r)=
 \left(w+\left\lfloor\frac{\bar r+1}{9}\right\rfloor,
       r+1\pmod9\right),\qquad
 \sigma_9^9(w,r)=(w+1,r).
\]
The second equality counts exactly one crossing from 8 to 0. Its first
coordinate records the advance that the second no longer distinguishes. In the
calendar of twelve windows, the exact sequence
\(0\to C_{12}\to C_{108}\to C_9\to0\) preserves carry between
windows. It does not split: \(C_{108}\) has exponent 108, and
\(C_{12}\times C_9\) has exponent 36. The proof and the carry law are
in \cref{prop:k-extension}.

Phase return thus acquires an operational meaning: it allows
readings made at the same position in the calendar to be compared while
the state preserves what happened between them. The holonomy \(\Gamma_9\)
is obtained from the enriched dynamics; \(K_{\rm ph}\) is its subsequent
reduced representation. The integer dodecaphasic register \(K\), which also uses
orientation and incidences, likewise has its own construction and
domain. Equality of a phase, a residue, or a expressed coordinate
does not by itself identify the histories. This is the initial information
preserved in passing to refinement.

\subsection*{Principle III. Compatibility of refinement and realization}
\hypertarget{prinh:refinamiento}{}

\emph{Refinement consists in adding resolution compatible with the prefix
while maintaining truncation relations and the masses of its prolongations.}
The reading at a finer scale specifies what occurred inside a
cylinder; forgetting that detail recovers the previous description.
The realization of a value arises when the cylinders of its reading
contract compatibly. The history determining those cylinders
remains in the domain of the construction.

In the tree of \cref{iv:cont:medida}, the initial set is finite and
nonempty, and each history has a finite, positive number of admissible
successors. A positive initial mass summing to one and local probabilities
\(p(\varepsilon\mid h)>0\), summing to one at each predecessor, produce
\[
 \mu(h\varepsilon)=\mu(h)p(\varepsilon\mid h),\qquad
 \sum_{\rho_k(y)=h}\mu(y)=\mu(h).
\]
Therefore, in the subsequent realization
\(H_k=L^2(X_k,\mu_k)\), the lift
\((J_kf)(y)=f(\rho_k(y))\) is isometric. Its adjoint is the conditional
expectation
\[
 (E_kg)(h)=\sum_{\rho_k(y)=h}
                  \frac{\mu(y)}{\mu(h)}g(y),
 \qquad E_kJ_k=I.
\]
Grouping the squared norm by predecessors proves the isometry, and grouping the
inner product proves the adjoint formula
(\cref{iv:cont:esperanza}). Lifting a reading and returning it
to the preceding level thus form a pair of explicit operations.

The hypotheses of finite branching and nonempty prolongation provide
compatible infinite histories. Their cylinders receive the preceding masses,
which determine a measure on the cylindrical sigma-algebra. The choice of
\(p\) belongs to the declared reader: the case of equiprobable successors gives
\(p=1/d(h)\), without imposing that measure on any other reading. This
construction preserves multiplicities of histories, including those that
share a realized value. Prolongability refers to the constructed
admissible tree; an additional restriction on routes must explicitly
preserve that property.

For the regional numerical reading, each block
\(b\in\mathbb F_3^6\) has an integer index
\(\nu(b)=\sum_{j=1}^6 b_j3^{6-j}\), lying between 0 and 728.
The reader fixes \(N_0=m\) and produces
\[
 N_{n+1}=729N_n+\nu(b_{n+1}),\qquad
 I_n=\left[\frac{N_n}{729^n},\frac{N_n+1}{729^n}\right].
\]
One has \(I_{n+1}\subseteq I_n\) and
\(\operatorname{diam}(I_n)=729^{-n}\). The left endpoints form
a rational Cauchy sequence; its class first defines the internal
realization of the reading. Subsequently identified with the usual
completion of the rationals, it corresponds to the unique point of the nested
closed intervals. The theorem following \eqref{eq:cilindro} proves
this statement and treats boundary conventions separately. For
regional characters, the effective reader of each decimal prefix
uses the computable brackets and the proved separation from
their boundaries, as specified in \cref{gen:prefijos-regionales}. This arbitrary
depth preserves the quantifier “for each finite length”; the
existence of a limit by compatibility and contraction is a distinct
statement and does not, by itself, provide an algorithm for every abstract
history. Histories representing the same point retain their
distinct records.

\subsection*{Principle IV. Correlative double reading and frame transport}
\hypertarget{prinh:doble}{}

\emph{The numerical and incidence readings arise from the same enriched
state and preserve their compatibilities under prolongation and changes of
frame.} The first specifies a realization through increasingly
narrow intervals. The second retains relations between positions, signs, and
boundaries. The holographic link consists in their common maps of origin and
transport, which allow the information expressed by each to be compared.

Let \(X_n^{\rm cal}\) be the domain of admissible prefixes including
their initial frame and successive frame transports. The arithmetic reading
uses the blocks of that prefix; the incidence reading uses the words
\(w_j\in\mathbb F_3^6\) and the frames \(f_j\). For a matrix
\(A\in M_6(\mathbb F_3)\), the linear reader
\[
 \gamma_A:\mathbb F_3^6\longrightarrow\mathbb F_3^{12},
 \qquad \gamma_A(w)=(w,wA)
\]
acts on words written as rows. Its specialization to the
Paley--Witt matrix constructed in \cref{exc:paley-codigo} produces the ternary
code and its incidences. The matrix is obtained in that finite
construction; a terminal real number does not select its entries. The order
\(w_6\to w_{12}\to w_{18}\to w_{24}\to w_{30}\to R_{36}\to G_9\)
maintains the transformations and memories used before each reading.

Compatibility can be checked on the operator itself. For
\(f\in\operatorname{GL}_6(\mathbb F_3)\), let
\(A^f=fAf^{-1}\). Then
\[
 \gamma_{A^f}(wf^{-1})
      =\gamma_A(w)(f^{-1}\oplus f^{-1}),
\]
since both sides are \((wf^{-1},wAf^{-1})\). In particular, the reading
of the prefix
\[
 Q_n(x)=\bigl(f_j,
       \gamma_{f_jA_Wf_j^{-1}}(w_j)\bigr)_{0\leq j<n}
\]
satisfies \(r_{n,m}Q_n=Q_m p_{n,m}\), where \(p_{n,m}\) truncates
histories and \(r_{n,m}\) truncates records. The causality of
frame updating guarantees that both procedures retain
exactly the same terms. The result
\cref{exc:limite-inverso} constructs its passage to the limit. Linear covariance
holds for the stated change of frame; additionally preserving a coordinate
support or a metric requires that structure to be transported, or
frame changes to be restricted to the group preserving it.

The correlative construction of the continuum extends this compatibility to
sheet transport, oriented balance, ternary incidence, the
rectangular port complex, and cellular memory with its fillings and
determinant lines. These aspects are constructed on the same
histories, not on five domains chosen afterward. The identities of
concatenation, naturality, and boundary ensure their assembly in
\cref{iv:rectangular-natural,iv:cont:memoria-natural,iv:cont:ensamblaje,iv:cont:mapa-correlativo}.
The five expository principles of this opening are not a new
partition of those five consubstantial constructions. Recovery
of a word by the first projection of \(\gamma_A\), for its
part, concerns that word: recovery of the complete enriched
history requires the other records distinguishing it to be preserved.

\subsection*{Principle V. Conservation and recovery of the complete record}
\hypertarget{prinh:conservacion}{}

\emph{Recoverable information is conserved jointly in the continuing
state, in the records produced, and in their correlations.} A
visible reading may contract while the complete record preserves an
isometric representation of the initial state. This principle allows
verifiable reversibility to be formulated: what is stored, in which
space it resides, and by which operator it is reconstructed are declared.

The transport of lifted histories in \cref{iv:cont:transporte} acts
between the spaces \(H_k=\ell^2(\widehat X_k)\) by
\[
 \mathcal U_ke_{\widehat h}=
 \sum_{\varepsilon\ \mathrm{admissible}}
 \sqrt{p(\varepsilon\mid\widehat h)}\,
 e_{\widehat h\varepsilon},\qquad
 \mathcal U_k^*\mathcal U_k=I.
\]
The second equality follows from local normalization and from distinct predecessors
having successors with distinct labels. If \(P_k\) reads sheet
\(\Sigma\) and \(Q_k=I-P_k\) reads sheet \(\Pi\), the operations
\[
 \begin{aligned}
 A_k&=P_{k+1}\mathcal U_kP_k+Q_{k+1}\mathcal U_kQ_k,\\
 \eta_k&=Q_{k+1}\mathcal U_kP_k+P_{k+1}\mathcal U_kQ_k
 \end{aligned}
 \qquad\text{satisfy}\qquad A_k^*A_k+\eta_k^*\eta_k=I.
\]
The Gram identity in \cref{iv:cont:gram} is proved by annihilating the
products of complementary projectors. It distinguishes remaining on a
sheet from changing sheets; the tritic value zero remains an
orientation and does not become a third sheet.

Write \(F_0=I\), \(F_{n+1}=A_nF_n\). Telescoping the
preceding identity gives, for every horizon \(N\),
\[
 \|x\|^2=\|F_Nx\|^2+
                 \sum_{n=0}^{N-1}\|\eta_nF_nx\|^2.
\]
The first term is the terminal residue that still continues; the others
are the recorded outputs. The theorem \cref{iv:cont:terminal} proves
that \(F_N^*F_N\) converges strongly to a positive operator
\(T_\infty\). One therefore obtains the complete record
\[
 \mathcal V:H_0\longrightarrow H_0\oplus
                     \bigoplus_{n\geq0}H_{n+1},\qquad
 \mathcal Vx=\bigl(T_\infty^{1/2}x,
                         (\eta_nF_nx)_{n\geq0}\bigr),
 \qquad\mathcal V^*\mathcal V=I.
\]
Indeed, the sum of squared norms is \(\|x\|^2\);
polarization gives the last identity. Its image \(\mathcal M\) is
closed, and the inverse on it is \(\mathcal V^*\). The terminal is
preserved until its extinction is proved; \(T_\infty\) need not be
a projector, and convergence of \(F_Nx\) is not assumed.

Conservation also extends to the operators reading that information.
For a bounded operator \(B\) on \(H_0\), the transport
\(B^{\mathcal V}=\mathcal VB\mathcal V^*|_{\mathcal M}\) preserves
products, adjoints, and inner products:
\[
 (B_1B_2)^{\mathcal V}=B_1^{\mathcal V}B_2^{\mathcal V},\qquad
 \langle\mathcal Vx,B^{\mathcal V}\mathcal Vy\rangle
                       =\langle x,By\rangle.
\]
Inserting \(\mathcal V^*\mathcal V=I\) proves both formulas. In the
nonadic reader of \cref{lib01:isometria}, the same identities are
realized with \(T_n=(8I+C_n)/9\) and
\(D_n=\sqrt8(C_n-I)/9\), for unitaries \(C_n\) on the declared
carrier. The composition retains its status as a reader and does not
automatically identify every \(C_n\) with the generating holonomy \(\Gamma_9\).
Recovery thus preserves the correlations of the record and the
transported observables within the precise image on which it acts.

The register \(K\) makes this distinction visible in a finite object. Its
encoding \(\kappa(K)\), in base \(B=1000\) and length twelve, is
\[
 \kappa(K)=\frac{N(K)}{B^{12}-1},\qquad
 N(K)=\sum_{j=1}^{12}K_jB^{12-j},\qquad
 K\in\{0,\ldots,999\}^{12}.
\]
Multiplying by \(B^{12}-1\) and performing twelve Euclidean divisions
recovers the register exactly, including its initial zero blocks.
The proposition \cref{lib01:inversion-kappa} further proves that an
absolute error less than \(1/[2(B^{12}-1)]\) allows the same integer to be recovered
by rounding. Reversibility concerns the register in its chart and
with its metadata; it does not identify the register with the entire TPK history. The
\emph{holographic arithmetic--geometric coupling constant} denotes,
in this article, the object \(K\) typed with its readers and relations. Its
representation arises from the enriched state and subsequently participates in
alpha closure: it is not used as a generating input for \(\pi\), nor
does its name make every scalar reading a universal invariant.

\subsection*{A joint reading of the five principles}

The sequence offers a positive criterion for navigating the article. Relational
unity fixes what is operated on; prolongation preserves what has been traversed;
refinement organizes resolution; double reading keeps value
and incidence together; recovery identifies the exact conditions for
reconstructing the record and its observables. The proofs are linked
because they act on objects of common provenance and respect their transport
maps. Wherever a reading forgets information, its domain and its
loss fiber remain defined.

The historical extreme and mean ratio shows how a relation between the whole
and its parts can determine a proportion. The holographic construction
developed here additionally studies relations conserved when resolution,
frame, and the distribution of information between continuation and
memory change. The self-scaling equation establishes a precise mathematical
link with that proportion; the identities of concatenation, refinement, and
isometry provide the additional content developed in this article.
The historical comparison thus retains its orienting function, while
each holographic conclusion rests on its own operations and proofs.

\subsection*{Primary sources for the historical comparison}

\begin{enumerate}[label=\arabic*.]
 \item Euclid, \emph{Elements}, Book I, Postulates 1--5. Greek
 text established by J. L. Heiberg, with a translation by Richard Fitzpatrick,
 corrected edition of 2008, p.~7; Book VI, Definition 2 in this edition,
 p.~156. \url{https://farside.ph.utexas.edu/books/Euclid/Elements.pdf}.
 \item Euclid, \emph{Elements}, digital edition by David E. Joyce,
 Clark University: Book I, Postulates 1--5,
 \url{https://mathcs.clarku.edu/~djoyce/elements/bookI/bookI.html};
 Book VI, Definition 3,
 \url{https://mathcs.clarku.edu/~djoyce/elements/bookVI/defVI3.html}.
 The difference in Book VI numbering is retained to facilitate
 comparison between editions.
\end{enumerate}
